\section{Reset minorization and central relative increments}

\label{special:reset}

The argument begins with a quantitative marginal theorem for every fixed positive omitted fraction, in both coordinate orders. That theorem produces an exact reset minorization. Two copies sharing a switch environment then have a central relative-permutation increment, whose return probability controls conditional uniformity.

The preceding route uses preparation domination to average a positive operator. Here we instead extract an exact uniform component from a long reset block. Uniformly placing the free labels makes the relative increment of two copies central, so a return probability can be computed by characters. Here the environment records every reset coin and the paths of the nonhole labels during the active sweeps.

Use the model and composition convention of Section~\ref{sec:conditional}. In this section $N=n=2^d$, and ``holes'' denotes the labels whose conditional assignment we will randomize.

\begin{theorem}[Reset construction]\label{thm:reset}
For $n=2^d$, the full-deck law from any deterministic start has total-variation distance tending to zero from uniform after $171798691968d$ physical shuffles.
\end{theorem}

\subsection{Partial mixing and forward minorization}

The reset needs an ordered marginal estimate for every fixed positive omitted fraction, rather than only for a half-deck. We prove it through a two-sweep signed-mass estimate, then use the same marginal bound in the reverse coordinate order to minorize the forward reset kernel.

Fix some observed labels and let $h$ be the number of remaining labels. Lemma~\ref{lem:slot-kernel} gives independent fair coins on the unobserved switches. To describe the conditional distribution for the position of a single particular hole label, one can evolve probability masses on the holes: at a pair of two holes the outgoing masses are both the average of the incoming ones, and at a mixed pair the hole mass goes to the outgoing hole. Which positions are holes at all the times is known from the observation. Although this description is conditional on the full observation, the forward calculation of mass through any layer only uses the observed paths up to then, as unobserved switches are left independent even when conditioning on the future part too. Uniform mass \(1/h\) per hole follows the same rule. We will bound the expected squared Euclidean norm of signed masses (mass zero off the holes, total mass zero) evolving by these rules, taking expectation when the trajectories are random, not conditioned. This is an adapted evolution in the switch process; at a mixed pair the placement of the outgoing mass uses the fresh fair switch coin. At every update the squared norm is non-increasing.

\begin{lemma}[Two-sweep mass contraction]\label{lem:reset-mass-contraction}
Fix $0<\rho\le1$. For all sufficiently large $d$ depending only on $\rho$, any deterministic hole placement with $h\ge\rho N$ and any initial signed masses supported on the holes and of total mass zero have expected squared norm after two sweeps at most
\[
 N^{-\rho/16}L,
\]
where $L$ is their initial squared norm. The masses evolve by the conditional averaging and transport rules just described. The estimate also holds conditionally on the past when the two sweeps begin.
\end{lemma}
\begin{proof}
Suppose \(h\ge\rho N\), and start at deterministic positions and signed masses with squared norm \(L\). A block of level \(i\) is a set with the last \(d-i\) bits fixed (the coordinates here are listed in the sweep order), of size \(2^i\). Write \(H_x\) for the hole indicator and \(M_x\) for the mass, and use block subscripts for sums over a block. Let \(r=\lfloor d/2\rfloor\). After one sweep use primes to denote the indicators, counts and masses then. Conditional on the first \(r\) layers, the updates along the last \(d-r\) coordinates proceed independently for different prefixes of length \(r\). Any fixed level-\(r\) block \(B\) has one vertex for each such prefix. Over the remaining layers the expectations per output of \(H_x, M_x\) evolve by ordinary pair averaging, and for \(M_x^2\) there is an upper bound by ordinary pair averaging (the square of a pair average is at most the average of the squares). Ordinary pair averaging through these layers computes the average within a prefix, successively averaging over each of the other coordinates. So \(H_B'\) is conditionally a sum of independent Bernoulli variables with mean sum \(h/2^{d-r}\), and the conditional mean of \(M_B'\) is zero. By the same independence, the conditional variance of \(M_B'\) is bounded by the sum of the output second moments in \(B\), at most \(L/2^{d-r}\) since the norm squared after the first \(r\) layers is at most \(L\). It follows (by the Chernoff bound for the counts) that the event \(D\) that all level-\(r\) blocks have hole density at least \(\rho/2\) after this sweep satisfies
\[
\Pr(D^c)\le 2^{d-r}\exp(-\rho 2^r/8),\qquad
\mathbb E\left[{\bf1}_D\sum_{B\ {\rm level}\ r}\frac{(M_B')^2}{H_B'}\right]\le \frac{2L}{\rho 2^r},
\]
where the indicated sum with the density requirement is only used on \(D\).

Here is an estimate for the next sweep from a deterministic input satisfying the density requirement (in the following calculation, input counts and masses refer to that input). Within a block \(B\) of level \(i\), consider the moments per output vertex after the first \(i\) layers; they are the same for each of those vertices, as follows recursively. In notation for a block, write \(u=H_B/|B|,\ m=M_B/|B|\); these are the first moments of indicator and mass. Write \(q\) for the second moment of mass. For the children of a parent use indices 1 and 2. Before the parent merge, the inputs to a switch come independently from the two child blocks (the layers so far act within the children with disjoint sets of coins). The new \(u,m\) are the averages of the two respective values, and
\[
q=\frac{q_1+q_2}{2}-\frac{q_1 u_2+q_2 u_1-2m_1m_2}{4}.
\]
In fact at a two-hole pair the per-output second moment loses the squared mass difference divided by 4 compared to just random placement; the correction averages that loss using independence (mass is zero off the holes). These calculations induct from the singleton blocks and in particular justify using position-independent output moments within the blocks.

At positive densities let \(a=m/u\) and \(v=q-m^2/u\ge 0\) (use Cauchy--Schwarz and support on the holes). For parents above level \(r\) all densities, including the children's, are at least \(\rho/2\). Rearrangement gives
\[
v=\tfrac12[(1-u_2/2)v_1+(1-u_1/2)v_2]+(1-u)J_B,\qquad
J_B=\tfrac12(u_1 a_1^2+u_2 a_2^2)-u a^2\ge 0.
\]
Weighting by block sizes, the sum of \(|B|v_B\) at level \(i>r\) is at most \(1-\rho/4\) times that at level \(i-1\), plus the sum of \(|B|J_B\) at level \(i\). The sum of \(|B|v_B\) at level \(r\) is bounded by the sum using second moments instead, at most the input squared norm. Each \(|B|J_B\) gives the sum for the two children minus the value for the parent of the quantity \(|B|u a^2=M_B^2/H_B\); the sum above level \(r\) telescopes and is bounded by \(\sum_{B\ {\rm level}\ r} M_B^2/H_B\). At the root \(v=q\) since the total mass is zero. Thus, iterating, for the second sweep on the event \(D\) the expected squared norm conditional on its input (which has squared norm at most \(L\)) is bounded by
\[
(1-\rho/4)^{d-r}L+\sum_{B\ {\rm level}\ r} (M_B')^2/H_B'.
\]
Using the pathwise bound \(L\) on \(D^c\), the expectation after the two sweeps is at most
\[
\left[(1-\rho/4)^{d-r}+\frac{2}{\rho 2^r}+2^{d-r}e^{-\rho2^r/8}\right]L\le N^{-\rho/16} L
\]
for \(d\) sufficiently large (the first term decays at least as fast as \(e^{-\rho d/8}\)). This estimate holds each time, conditionally on the past, with arbitrary input meeting just \(h\ge\rho N\) and the signed-mass conditions (the density event concerns the intermediate time).
\end{proof}

\begin{proposition}[Partial mixing at fixed omitted fraction]\label{prop:reset-partial-mixing}
For fixed $0<\rho\le1$, set $s=\lceil320/\rho\rceil$. For all sufficiently large $d$ depending only on $\rho$, after $2s$ sweeps the positions of any ordered tuple of $b\le(1-\rho)N$ distinct labels are within $\varepsilon_N=N^{-8}$ in total variation of the uniform ordered injection, from any deterministic start. The same assertion holds for sweeps in reverse order $d,\ldots,1$.
\end{proposition}
\begin{proof}
By Lemma~\ref{lem:reset-mass-contraction}, after $s$ pairs of sweeps the bound is $N^{-20}L$, since $s\rho/16\ge20$.

Take the ordered tuple in the proposition, indexing its labels by \(1,\ldots,b\). For label \(i\), condition on the trajectories of labels 1 through \(i-1\). The number of holes is at least \(\rho N\). Initialize signed mass by the point mass for \(i\) minus uniform mass on the holes, so its squared norm is at most 1. By the conditional trajectory description the final signed mass is the endpoint probability difference from uniform on positions other than the previous labels' endpoints. Its expected total variation norm is at most \(\frac12\sqrt{N}\, N^{-10}\) by Cauchy--Schwarz and the squared-norm bound. Conditioning on just the previous labels' endpoints instead does not increase this expected bound, by averaging (the uniform reference is unchanged given those endpoints). Total variation from the uniform joint injection is bounded by summing these conditional variation estimates: telescope the laws using the actual distribution for an initial segment and the uniform sequential conditional kernels for the remaining labels. This proves the asserted total-variation bound. Its proof only used the order of the coordinates as an indexing, so works in the reversed order as well.

\end{proof}

\begin{corollary}[Forward reset minorization]\label{prop:reset-minorization}
Let \(K\) be the kernel of the \(2s\) forward sweeps in Proposition~\ref{prop:reset-partial-mixing} on an ordered injection state space \(\Omega\) with \(b\le (1-\rho)N\) labels. Then
\[
K^2(x,y)\ge \frac{p_0}{|\Omega|},\qquad p_0=1-2\varepsilon_N.
\]
\end{corollary}
\begin{proof}
The kernel $K$ is doubly stochastic. Its transpose $K^*$ uses the inverse permutation law, sweeps in reversed order, so Proposition~\ref{prop:reset-partial-mixing} applies row-wise to both kernels. Write
\[
K^2(x,y)=\sum_z K(x,z)K^*(y,z).
\]
In this inner product truncate each factor above by \(w=1/|\Omega|\) (take the minimum). The truncated factors are \(w-a_z,w-b_z\) with nonnegative deficits satisfying \(\sum_z a_z\le\varepsilon_N\) and \(\sum_z b_z\le\varepsilon_N\); their product is at least \(w^2-w(a_z+b_z)\), giving the displayed lower bound.

\end{proof}

\subsection{Comparison conditional on a switch environment}

We have proved uniform marginal control in both coordinate orders and converted it into pointwise minorization. The reverse-order estimate was used to control columns of the forward kernel; the actual reset is still a product of forward sweeps. We next exploit the minorization without changing that actual run.

 From now on let \(\rho=2^{-20}\), with \(d\) sufficiently large, and fix a set of \(k=\rho N\) nonhole labels; let \(H\) be the set of \(h=N-k\) hole labels. These roles of the labels are for the comparison only. Run \(l=128\) blocks, each consisting of a reset of \(4s\) sweeps followed by a single mixing sweep (all sweeps are actual random shuffles). Write \(X\) for the process as permutations taking labels to positions.

Define an environment by revealing all the reset coins and, during the single mixing sweep of each block, revealing the switches touched by the nonhole trajectories and their values. The environment determines all nonhole paths. Conditional on it, the unrevealed switches during mixing sweeps (the two-hole switches) are independent and fair, by the fixed-values observation as for trajectories earlier: every completion has the same nonhole paths, using also the fixed reset values. We may generate \(Y\) independently of \(X\) conditional on the environment with the same conditional law, by sharing the reset and touched switch values and independently redrawing the two-hole switches during mixing sweeps. The copies have the same nonholes in the same positions. Unconditionally this joint generation can be done chronologically, generating fresh coins for \(X\) and using independent coins for the redraws (the switches to be redrawn in a layer are known upon entering it).

For this environment $\mathcal E$, let $p_{\mathcal E}$ be the final law of the hole assignment conditional on $\mathcal E$, a bijection from the $h$ hole labels to their available endpoint positions, and let $u_{\mathcal E}$ be uniform on those $h!$ assignments.

\begin{proposition}[Uniformity conditional on the reset environment]\label{prop:reset-conditional}
For the run just defined, with $\rho=2^{-20}$, $s=\lceil320/\rho\rceil$ and $l=128$ blocks of $4s$ reset sweeps followed by one mixing sweep,
\[
 \mathbb E\|p_{\mathcal E}-u_{\mathcal E}\|_{\rm TV}\longrightarrow0
 \qquad(d\longrightarrow\infty),
\]
uniformly over deterministic starting decks. The environment $\mathcal E$ reveals every reset coin and the nonhole paths during the mixing sweeps, as above.
\end{proposition}
\begin{proof}
The relative permutation \(X^{-1}Y\) acts on \(H\), starts at the identity and does not change during a reset. In a layer of a mixing sweep it is multiplied on the left by a product of relative transpositions: transpose the two incident labels of \(X\) at a two-hole switch if the value for \(Y\) differs. Indeed if the switches act by \(T_X,T_Y\) on positions and \(X,Y\) here denote the pre-layer values, the new relative permutation is \((X^{-1}T_X^{-1}T_Y X)X^{-1}Y\). Given the \(X\) switches and paths, these relative toggles are independently fair. We write \(Q\) for the resulting multiplier over one mixing sweep; it only uses the hole injection for \(X\) at the sweep start and the switches (including redraws) in that sweep, not the prior relative permutation.

For now mark two types of events for \(X\) only (and these marks do not change its run). At each reset mark an acceptance such that the injection of the \(h\) ordered hole labels just afterwards, jointly with acceptance, has probability \(p_0/|\Omega|\) for each injection conditional on the past. This is the elementary minorization-splitting mechanism developed systematically by Nummelin~\cite[Section~2]{nummelin1978}. Here the calculation is finite and explicit. By the \(K^2\) bound applied with \(b=h\), given the start \(x\) and sampled endpoint \(y\) of the hole labels, accept with probability \(p_0/(|\Omega|K^2(x,y))\) using independent extra randomness. The statement holds also with the past including the other copy in the chronological construction (we do not condition here on the environment, and the relative permutation is unchanged during the reset).

To describe the second mark, consider a mixing sweep with hole injection for \(X\) uniform at its start, as upon acceptance. The graph on all positions at the start, connecting two whose paths share a switch in \(X\), is simple and \(d\)-regular: after two paths meet they differ in that updated bit which doesn't change again during the sweep, so they cannot meet at a switch along any later coordinate. This graph is defined by the fresh \(X\) coins independent of the hole injection; given the graph the nonholes are on a uniform \(k\)-subset of vertices. The probability any vertex has at least \(\lceil d/2\rceil\) nonhole neighbors is at most
\[
N2^d\rho^{d/2}=N^{-8}
\]
by a union bound and sampling without replacement. Mark the mixing sweep as good for \(X\) on the complementary event; every hole label then has at least \(d/2\) two-hole switches on its path.

Let \(\nu\) on the symmetric group on \(H\) be the subprobability law of \(Q\) with this goodness requirement when the \(X\) hole injection starts uniformly. Its mass \(q\) is at least \(1-N^{-8}\). The increment and goodness depend on no nonhole identities, and this law is central (invariant under conjugation) by the symmetry of the uniform injection under hole relabelings. After a reset, jointly requiring acceptance and goodness gives exactly this sublaw with the additional factor \(p_0\): writing \(\mathcal F\) for the past before the reset in the chronological construction, we have
\[
\Pr(\text{accept},\text{ good},Q=g\mid \mathcal F)=p_0\nu(g).
\]
In particular if \(G_X\) denotes all accepts and all good marks for \(X\), and \(S\) denotes the final relative permutation, we have
\[
\Pr(G_X)=(p_0 q)^l,\qquad
\Pr(S=\mathrm{id},G_X)=p_0^l\,\nu^{*l}(\mathrm{id}).
\]
Indeed, for the relative multiplication the placement of \(X\)'s holes after the accepted reset is uniform independent of the past, regardless of the current relative permutation. The relative coin differences in the next sweep come from the independent redraws and give the multiplier law described above, allowing iteration with convolution.

We record two properties used to estimate this convolution. Conditional on \(X\) in a good mixing sweep, the sign of \(Q\) is fair since there is a two-hole switch and the relative toggles are independent. And for any \(A\subset H\),
\[
\nu\{Q:\ Q\text{ fixes } A \text{ pointwise}\}\le 2^{-d|A|/4}\le h^{-|A|/4}.
\]
For this estimate, imagine using the switches for \(Y\) but starting the mixing sweep at the very same permutation as \(X\). The final relative permutation for this sweep would be \(Q\). Within one sweep a path with given start and end is unique (updated bits equal end bits and not-yet-updated bits equal start bits). Fixation of a label thus requires agreement at every switch on its \(X\) path. In a good sweep at least \(d|A|/4\) distinct two-hole switch coins are involved on the paths for \(A\), counting any such switch at most twice towards the path incidences. Their agreement probability gives the bound.

\subsection{Return estimate for the central sublaw}

The accepted reset has made the relative increments central and independent of their previous product. The remaining question is how quickly their convolution puts mass $1/h!$ at the identity. This return estimate will control the mean squared density of a conditional assignment, rather than an unconditional marginal. We show
\[
h!\,\nu^{*l}(\mathrm{id})\le q^l+o(1).
\]
Use Lemmas~\ref{lem:character-fixed-points} and~\ref{lem:degree-sum} with group size parameter $h$. For a permutation $\sigma$ of $H$, let $F(\sigma)$ count its fixed points. Write \(\theta_\lambda=\sum_\sigma\nu(\sigma)\chi_\lambda(\sigma)/d_\lambda\), where $d_\lambda$ is the irreducible degree. Taking $\delta=1/32$ in Lemma~\ref{lem:character-fixed-points}, for \(F\le h^{7/8}\) the normalized absolute character is at most \(d_\lambda^{-1/32}\) for large \(h\). For \(f>h^{7/8}\) the pointwise fixation estimate gives
\[
\nu(F=f)\le \binom h f h^{-f/4}\le h^{-f/10}
\]
for large \(h\), using \(\binom h f\le(eh/f)^f\). The all-fixed-point bound~\eqref{eq:character-all-fixed} gives
\[
|\chi_\lambda(\sigma)|/d_\lambda\le d_\lambda^{-1/32} h^{f/32}.
\]
The tail satisfies
\[
 \sum_{f>h^{7/8}}h^{-f(1/10-1/32)}\le 1
\]
for large $h$. Summing these character estimates against \(\nu\) therefore gives \(|\theta_\lambda|\le 2 d_\lambda^{-1/32}\).

By centrality each averaged representation matrix is scalar by Schur's lemma, with scalar \(\theta_\lambda\). Taking the regular representation trace for the convolution therefore gives \(h!\,\nu^{*l}(\mathrm{id})=\sum_\lambda d_\lambda^2 \theta_\lambda^l\). The trivial term is \(q^l\) and the sign term is zero. With \(l=128\) the sum of absolute values of the other terms is at most \(2^l\sum_\lambda d_\lambda^{-2}\) excluding those two terms, hence is \(o(1)\) by Lemma~\ref{lem:degree-sum}. This proves the return estimate.

\subsection{From the return estimate to conditional uniformity} Write \(p_{\mathcal E}^G\) for the subprobabilities of \(p_{\mathcal E}\) including the requirement \(G_X\). For the conditionally independent copy \(Y\), define its own event \(G_Y\) by the same acceptance and goodness requirements, with independent auxiliary randomness. By the triangle inequality with the subprobabilities and Cauchy--Schwarz,
\[
\mathbb E\|p_{\mathcal E}-u_{\mathcal E}\|_{\rm TV}
\le \tfrac12\Pr(G_X^c)
+\tfrac12\sqrt{h!\Pr(S=\mathrm{id},G_X,G_Y)-2\Pr(G_X)+1}.
\]
Indeed \(p_{\mathcal E}-p_{\mathcal E}^G\) is nonnegative of expected mass \(\Pr(G_X^c)\), and for the remaining difference with uniform the expression inside the square root is the expectation of \(h!\sum(p_{\mathcal E}^G-1/h!)^2\). This identity uses the two conditionally independent copies and their separate marks; assignment agreement is \(S=\mathrm{id}\). Dropping \(G_Y\) from the positive term and using the convolution formulas with the return estimate bounds the expression by \(1-(p_0q)^l+o(1)=o(1)\); likewise \(\Pr(G_X^c)=o(1)\). This truncation by good subprobabilities only costs small mass in total variation. All estimates are uniform over the deterministic starting deck, proving Proposition~\ref{prop:reset-conditional}.
\end{proof}

\subsection{Full-deck mixing}
\begin{proof}[Proof of Theorem~\ref{thm:reset}]
\label{special:reset:bound}
Unconditionally, the nonhole marginal is close to uniform: the first \(2s\) sweeps already have the partial mixing estimate for its \(k\) labels, and subsequent independent switches in the actual run preserve the uniform marginal kernel-wise, so they cannot increase its distance from uniform. Thus its total variation from uniform is at most \(\varepsilon_N\). The distribution obtained by assigning holes uniformly given the environment has these very nonhole endpoints and uniform assignment on their complement, so is within \(\varepsilon_N\) of uniform on the whole deck. By the conditional bound, the actual endpoint distribution differs from this distribution by \(o(1)\) in total variation. The bound is independent of the initial deck. We have proved an upper bound on the required time for large \(d\) by \(128(4s+1)d\) physical shuffle steps.
With $\rho=2^{-20}$, $s=320\cdot2^{20}=335544320$, and the resulting upper time is $171798691968d$ complete shuffles.

\end{proof}

The support obstruction was proved in Section~\ref{sec:conditional}; in particular it implies the lower bound $t_{\mathrm{mix}}\ge d$ for all sufficiently large $d$.
